\subsection{Desarrollo teórico unificado}

La dificultad central del transporte cooperativo consiste en hacer compatibles varias
abstracciones que operan a escalas distintas. La asignación es discreta porque un AMR se
compromete con una carga concreta. La preferencia puede ser continua porque una dinámica
poblacional reparte masa entre alternativas. La factibilidad mecánica es convexa a nivel
de esfuerzos admisibles, pero combinatoria a nivel de selección de contactos. La
seguridad se expresa como desigualdades sobre
trayectorias. La arquitectura propuesta se defiende precisamente porque conserva esos
niveles separados hasta el momento en que deben acoplarse.

Un subproblema elemental puede escribirse como
\[
  \min_{a,z,u,\lambda}
  \quad
  J_{\rm op}(a,z,u)
  + \sum_{k}\alpha_k\pos{D_k-S_k(C_k)}
  + \sum_{k}\beta_k\rho_k(C_k)
\]
sujeto a
\[
  a_i\in\mathcal L,\qquad
  z_{ikh}\in\{0,1\},\qquad
  u_i(t)\in\mathcal U_i,\qquad
  h_{io}(q_i(t))\geq 0,
\]
\[
  C_k=\{i:a_i=k\},\qquad
  0\leq\lambda\leq\bar f,\qquad
  \rho_k(C_k)=
  \min_{0\leq\lambda\leq\bar f}
  \norm{G_{C_k}\lambda-w_k^{\rm dem}}_2 .
\]
La formulación no se resuelve de forma monolítica en la memoria. Su función es fijar qué
debe conservar cualquier aproximación: demanda no servida, geometría de contacto,
saturación, seguridad, energía y coste de cómputo. Smith-QR actúa como una relajación
estructurada de este problema. La dinámica de Smith aproxima la presión de asignación,
el cierre QR produce una decisión entera, y las guardias físicas rechazan coaliciones que
no cumplen capacidad, wrench o seguridad.

\begin{definicion}[Juego físico de coalición]
Un juego físico de coalición para transporte AMR se define por
\[
  \mathcal G_{\rm phys} =
  \langle \mathcal R,\mathcal L,\{\mathcal A_i\},\{J_i\},
  \{F_k\},G(t)\rangle ,
\]
donde $\mathcal A_i$ contiene cargas y roles admisibles, $J_i$ es la utilidad local del
robot, $F_k$ agrupa las restricciones físicas de la carga $k$ y $G(t)$ limita la
información que puede circular. Una coalición $C_k$ es válida si satisface al mismo
tiempo cierre entero, capacidad efectiva, residual wrench, margen de seguridad y
retornabilidad energética.
\end{definicion}

Esta definición evita una confusión frecuente. Una partición de robots puede ser una
coalición en sentido combinatorio y, aun así, no ser una coalición física. La diferencia
no es terminológica: aparece en los experimentos cuando una asignación escalar completa
una demanda aparente pero no genera torque, cuando la carga llega con mala orientación o
cuando un robot asignado queda fuera de rango y deja de aportar información útil.

\begin{proposicion}[Monotonía del déficit de capacidad]
Sea $S_k(C)=\sum_{i\in C}c_i\eta_{ik}$ con $c_i\geq 0$ y
$\eta_{ik}\in[0,1]$. Si $C\subseteq C'$, entonces
\[
  \pos{D_k-S_k(C')}\leq \pos{D_k-S_k(C)} .
\]
\end{proposicion}

\begin{proof}[Demostración]
Como todos los términos añadidos son no negativos,
$S_k(C')=S_k(C)+\sum_{i\in C'\setminus C}c_i\eta_{ik}\geq S_k(C)$. La función
$x\mapsto\pos{D_k-x}$ es no creciente. Al componer ambas relaciones se obtiene la
desigualdad.
\end{proof}

La proposición justifica el uso de precios de déficit: añadir capacidad efectiva nunca
empeora el déficit de una carga. No dice que añadir robots sea siempre deseable. El coste
de viaje, la congestión, el uso de batería y la pérdida de disponibilidad para otras
cargas pueden hacer que una coalición sobredimensionada sea mala aunque reduzca déficit.
Por eso el payoff incluye penalización de exceso y distancia, no solo recompensa por
unirse a la carga más demandante.

\begin{proposicion}[Residual wrench como distancia a un conjunto admisible]
Para una coalición $C$, defínase
\[
  \mathcal W(C)=\{G_C\lambda:\ 0\leq\lambda\leq\bar f_C\}.
\]
Entonces $\rho(C)=\min_{w\in\mathcal W(C)}\norm{w-w^{\rm dem}}_2$ es la distancia
euclídea desde el wrench demandado hasta el conjunto de wrench admisibles. Si
$C\subseteq C'$ y los contactos añadidos no eliminan contactos previos, entonces
\[
  \rho(C')\leq \rho(C).
\]
\end{proposicion}

\begin{proof}[Demostración]
El conjunto $\mathcal W(C)$ es la imagen lineal de una caja cerrada y acotada, por lo que
es compacto. La distancia euclídea al conjunto alcanza mínimo. Si $C\subseteq C'$, toda
solución factible de $C$ también puede representarse en $C'$ asignando esfuerzo cero a
los contactos añadidos. Por tanto, $\mathcal W(C)\subseteq\mathcal W(C')$. Minimizar la
misma distancia sobre un conjunto mayor no puede aumentar el valor óptimo.
\end{proof}

Esta monotonía es más débil de lo que parece. El residual no aumenta al añadir un contacto
ideal, pero el sistema real sí puede empeorar por saturación, mala orientación, bloqueo
geométrico o coste de coordinación. La evaluación debe considerar $\rho(C)$,
$s_k^{\rm sat}$, el despeje de la carga y el tiempo requerido para ocupar el slot. SP3
mide la capacidad de producir wrench. SP4 exige una trayectoria segura. SP5 comprueba
la pose de la carga extendida.

\begin{lema}[Cierre entero bajo demanda visible]
Supóngase que todas las cargas son visibles, que $\sum_k n_k\leq N$, que cada robot puede
ser asignado a una sola carga y que el redondeo escoge robots sin repetir hasta cubrir las
demandas $n_k$. Si ninguna guardia física rechaza asignaciones, el cierre QR produce
$|C_k|\geq n_k$ para toda carga demandada.
\end{lema}

\begin{proof}[Demostración]
La hipótesis $\sum_k n_k\leq N$ implica que existe al menos una asignación entera con
quórum completo. El procedimiento de cierre consume a lo sumo un robot por selección y
no repite robots. Al iterar sobre las cargas demandadas, cada carga recibe robots hasta
alcanzar $n_k$ antes de pasar a capacidad sobrante. Como no hay rechazos físicos ni
invisibilidad, ninguna selección válida se descarta. Por inducción sobre las cargas, el
quórum queda cerrado en todas ellas.
\end{proof}

El lema delimita el papel exacto de QR. QR no prueba optimalidad global ni factibilidad
mecánica; prueba que la preferencia continua de Smith puede transformarse en un conjunto
entero cuando el problema discreto es materialmente cerrable. Si la demanda excede la
flota, si una carga no ha sido descubierta o si una guardia física descarta contactos, el
resultado correcto consiste en reportar déficit, reasignar o declarar inviabilidad.

\begin{proposicion}[Invariancia de seguridad bajo filtro de barrera]
Sea $\mathcal S=\{q:h(q)\geq0\}$ un conjunto seguro y supóngase que la acción aplicada
cumple
\[
  \dot h(q,u)+\alpha h(q)\geq0,\qquad \alpha>0 .
\]
Si $q(0)\in\mathcal S$, entonces la trayectoria continua permanece en $\mathcal S$.
\end{proposicion}

\begin{proof}[Demostración]
La desigualdad implica $\dot h(t)\geq-\alpha h(t)$. Por comparación escalar,
$h(t)\geq h(0)e^{-\alpha t}$. Como $h(0)\geq0$, se obtiene $h(t)\geq0$ para todo
$t\geq0$.
\end{proof}

En simulación discreta esta garantía se usa con prudencia. El paso temporal, la saturación
del uniciclo, la geometría rectangular de la carga y la presencia de varios robots pueden
introducir errores de discretización. Por eso la memoria no declara seguridad solo porque
un campo tenga término repulsivo; la declara cuando la trayectoria renderizada y las
métricas discretas conservan margen frente a obstáculo, sin colisión ni bloqueo persistente.

\subsubsection{Reparto cuadrático, coste de anarquía y cierre distribuido}

La capa de esfuerzo admite una solución cerrada que permite separar salud actuadora,
beneficio colectivo y coste individual. Sea $f_i\in\RR^2$ la fuerza del robot $i$,
$\eta_i>0$ su salud efectiva, $M_T$ la masa total y $s$ la aceleración de carga que debe
compensarse. Cada robot minimiza
\[
 J_i(f_i;f_{-i})=
 \frac{1}{2\eta_i}\norm{f_i}^2+
 \frac{\kappa}{2}\norm{s-\frac{1}{M_T}\sum_j f_j}^2.
\]
Defínanse $H=\sum_j\eta_j$ y $\gamma=\kappa H/M_T^2$.

\begin{proposicion}[Equilibrio del reparto cuadrático]
El juego anterior tiene un único equilibrio de Nash. La fuerza y el error común en ese
equilibrio son
\[
 f_i^{\rm NE}=\frac{\kappa\eta_i M_T}{M_T^2+\kappa H}s,
 \qquad
 e^{\rm NE}=\frac{1}{1+\gamma}s.
\]
\end{proposicion}

\begin{proof}[Demostración]
La condición de primer orden es
$f_i/\eta_i-(\kappa/M_T)e=0$, donde
$e=s-M_T^{-1}\sum_j f_j$. Al sumar las mejores respuestas se obtiene
$(1+\gamma)e=s$. La sustitución produce la expresión de $f_i^{\rm NE}$. La matriz
cuadrática es definida positiva en la decisión propia, lo que da unicidad.
\end{proof}

El equilibrio reparte el esfuerzo en proporción a $\eta_i$, pero cada robot internaliza
solo su coste y una fracción del beneficio común. El cociente entre el coste social en
equilibrio y su mínimo permite cuantificar esa pérdida.

\begin{proposicion}[Coste de anarquía del juego cuadrático]
Para $N$ robots, el coste de anarquía es
\[
 {\rm PoA}(\gamma,N)=
 \frac{(\gamma+N)(1+N\gamma)}{N(1+\gamma)^2},
 \qquad
 \max_{\gamma>0}{\rm PoA}(\gamma,N)=\frac{(N+1)^2}{4N}.
\]
El máximo se alcanza en $\gamma=1$.
\end{proposicion}

\begin{proof}[Demostración]
Para una fuerza total $F=\sum_i f_i$, el reparto de mínimo esfuerzo ponderado es
$f_i=(\eta_i/H)F$. El óptimo social minimiza
$\norm{F}^2/(2H)+(N\kappa/2)\norm{s-F/M_T}^2$ y produce
$F^{\rm SO}/M_T=N\gamma s/(1+N\gamma)$. El equilibrio produce
$F^{\rm NE}/M_T=\gamma s/(1+\gamma)$. Evaluar ambos costes y formar el cociente da la
expresión anterior. Su derivada se anula en $\gamma=1$.
\end{proof}

Cuando el juego incorpora una restricción compartida de fuerza y torque, el equilibrio
variacional recupera el reparto ponderado de mínimo esfuerzo. Para la componente
traslacional, la participación de cada robot es $\eta_i/H$. La componente rotacional
añade el brazo $(Rr_i)^\perp$ y el momento ponderado
$S=\sum_j\eta_j\norm{r_j}^2$. Esta relación conecta el reparto estratégico con la
geometría de contacto usada en SP3.

La ejecución distribuida reemplaza los agregados globales por estimaciones locales. En
el modelo linealizado, sea $\mu$ la constante de monotonicidad fuerte, $\vartheta$ el
acoplamiento transversal, $c$ la ganancia de consenso y $\lambda_2(L_G)$ la conectividad
algebraica. La condición suficiente
\[
 c\lambda_2(L_G)\mu>\vartheta^2
\]
hace que la contracción del juego y del consenso domine el acoplamiento cruzado. Esta
desigualdad se aplica al buscador simplificado y no constituye una garantía para contacto
tridimensional, saturación o pérdida persistente de conectividad. La frontera anterior se conserva como resultado del buscador simplificado, pero deja de usarse como evidencia numérica por identidad. La nueva V1 resuelve de forma independiente el sistema KKT del reparto de wrench y lo compara con la ley cerrada (error máximo $2{,}85\times10^{-14}$). V2 contrasta la identidad de potencia Lagrange--Hamilton mediante diferencias finitas y la proyección HOCBF contra un QP SLSQP independiente. V3 verifica el certificado ISS práctico que sigue.

\begin{teorema}[ISS práctica común frente al error de realización de wrench]
\label{thm:iss-practica-coalicion}
Sea $z=(e,\dot e)$ el error de una carga fija controlada por realimentación PD,
\[
 \dot z=Az+B d_{\sigma(t)},
\]
donde $A$ es Hurwitz, el cambio de coalición no reinicia el estado y $\lVert d_{\sigma(t)}\rVert\leq\bar d$. Si $P\succ0$ satisface $A^\top P+PA=-I$, entonces, para cualquier señal medible de cambio $\sigma(t)$ y cualquier $0<\varepsilon<1$,
\[
 \dot V\leq-\alpha V+\sigma_d\bar d^2,
 \qquad V=z^\top Pz,
\]
con
\[
 \alpha=\frac{1-\varepsilon}{\lambda_{\max}(P)},
 \qquad
 \sigma_d=\frac{\lVert PB\rVert_2^2}{\varepsilon}.
\]
En consecuencia,
\[
 \lVert z(t)\rVert
 \leq
 \sqrt{\frac{V(0)e^{-\alpha t}+(\sigma_d/\alpha)\bar d^2(1-e^{-\alpha t})}{\lambda_{\min}(P)}}.
\]
\end{teorema}

\begin{proof}[Demostración]
La ecuación de Lyapunov da $\dot V=-\lVert z\rVert^2+2z^\top PBd_{\sigma(t)}$. La desigualdad de Young produce $2z^\top PBd\leq\varepsilon\lVert z\rVert^2+\lVert PB\rVert_2^2\lVert d\rVert^2/\varepsilon$. Como $V\leq\lambda_{\max}(P)\lVert z\rVert^2$, se obtiene la desigualdad diferencial. El lema de comparación y $V\geq\lambda_{\min}(P)\lVert z\rVert^2$ proporcionan la cota.
\end{proof}

La matriz cerrada es común porque la masa, la inercia y las ganancias de la carga no cambian durante la sustitución; la coalición modifica el error de realización $d_{\sigma(t)}$. Por ello no se impone tiempo de permanencia. Si cambia $A$, existe un reset de estado o el error deja de estar acotado, el teorema no aplica y sería necesario un argumento de Lyapunov múltiple. Para los parámetros de la campaña, V3 obtuvo $\lambda_{\min}(P)=0{,}278$, residual de Lyapunov $9{,}04\times10^{-16}$ y ninguna violación de la cota en las trayectorias lineales independientes.
\subsubsection{Acoplamiento entre potencial, mercado y ejecución}

La lectura de potencial permite unir asignación y control sin confundirlos. Sea
\[
  \Phi(a)=
  -\sum_k\alpha_k\pos{D_k-S_k(C_k)}
  -\sum_k\beta_k\rho_k(C_k)
  -\sum_i\gamma_i d(q_i,q_{a_i})
  -\sum_i\zeta_i\pos{b_{\min}-b_i}.
\]
Si el payoff local aproxima la diferencia marginal
\[
  J_i(k;a_{-i})\approx
  \Phi(k,a_{-i})-\Phi(\varnothing,a_{-i}),
\]
entonces cada robot recibe una señal coherente con el objetivo común. Entrar en una carga
solo tiene sentido si reduce déficit, residual o coste operativo de forma suficiente. El
precio de déficit $\pi_k$ puede verse como multiplicador local de la restricción
$S_k(C_k)\geq D_k$:
\[
  \pi_k(t+1)=\pos{\pi_k(t)+\eta_\pi\bigl(D_k-S_k(C_k,t)\bigr)}.
\]
Cuando falta capacidad, $\pi_k$ sube y atrae robots; cuando hay exceso, la presión cae.
La misma idea se aplica al residual wrench con un precio mecánico $\mu_k$:
\[
  \mu_k(t+1)=\pos{\mu_k(t)+\eta_\mu\bigl(\rho_k(C_k,t)-\epsilon_k\bigr)}.
\]
El payoff útil contabiliza escasez, distancia, batería y viabilidad mecánica mediante
señales locales.

Esta construcción explica la comparación entre métodos heterogéneos. Un greedy
cercano prueba si basta con distancia. Un Hungarian centralizado prueba si basta con
optimización escalar. CBBA prueba si una subasta local con consenso es suficiente. Las
variantes aprendidas prueban si una política entrenada captura interacciones no lineales.
Los oráculos fijan cotas. Smith-QR y sus extensiones prueban la hipótesis específica de
la memoria: una señal local con precios físicos y cierre entero puede competir mejor en
los regímenes donde la tarea deja de ser una unidad indivisible y pasa a ser una carga
compartida.

\subsubsection{Correspondencia entre afirmación y evidencia}

La teoría anterior conduce a un criterio de evaluación severo. Una afirmación del tipo
``el método asigna bien'' queda incompleta si no especifica qué nivel superó. En esta
memoria, una solución puede fallar por ocho razones distintas: no cierra quórum, cubre
cardinalidad pero no capacidad, cubre capacidad pero no wrench, genera wrench pero no
trayectoria, llega con robots individuales pero no con la carga, no se recupera ante
fallo, depende de comunicación global o deja de ser practicable al crecer la flota.

La escalera SP1--SP8 impide que una mejora local se confunda con una solución general.
Partiendo del reclutamiento homogéneo unificado SP1, cada subproblema añade una restricción que
vuelve falso algún éxito anterior. Si SP1 se supera, todavía puede fallar SP2; si SP2 se
supera, todavía puede fallar SP3; si SP3 se supera, todavía queda movimiento seguro,
pose, recuperación, comunicación y escala. La contribución queda formulada por capas y
se identifica dónde la arquitectura distribuida conserva valor frente a familias
clásicas, aprendidas y centralizadas.
